\sect{14. Orientación y pistas}
\textbf{Tu misión, paso a paso.} Explora cada planta, busca una ruta hacia las zonas inferiores, reúne las llaves y piezas necesarias y conserva un grupo capaz de combatir y curarse. El objetivo final es encontrar y derrotar a Xanathar. Las habitaciones secretas y los encargos opcionales pueden recompensar la exploración, pero encontrar todos los secretos no es lo mismo que cumplir la misión principal.

\textbf{Automapa.} La aventura aparece en la salida del V9990 y el mapa en la del MSX. Conecta dos pantallas o cambia de entrada en el monitor. En openMSX, elige la fuente de vídeo MSX para ver el mapa y GFX9000 para volver a la aventura.

El norte queda arriba. La flecha blanca indica dónde está el grupo y hacia dónde mira. Las casillas negras están sin explorar. El mapa marca las puertas visibles en naranja, las escaleras de subida en cian y las de bajada en amarillo. Vuelve a un lugar para ver si algo ha cambiado. Atraviesa las paredes ilusorias para que aparezcan como pasadizos.

El mapa se guarda con el progreso. Abre \texttt{CAMP} para detener la aventura mientras lo consultas.

\textbf{Consultar otras plantas.} \key{F2} muestra la planta anterior, \key{F3} la siguiente y \key{F4} vuelve a la planta en la que estás. También funciona en el campamento. Los huecos abiertos visibles se marcan en rojo y los portales en magenta. Anota por separado las inscripciones y los destinos de los portales.

\subsection*{Primeros minutos}
\begin{enumerate}
\item Coloca a los mejor protegidos delante. Comprueba que llevan sus armas en la mano, no en la mochila.
\item Abre los inventarios, identifica libro, símbolo, comida y munición. Lee la carta de encargo.
\item Prepara curación y algunos ataques mágicos; descansa en un lugar seguro para recuperarlos.
\item Explora un tramo corto, observa brújula y paredes, recoge lo útil y prueba los botones visibles.
\item Guarda antes de bajar por primera vez. Conserva una ruta de regreso a un lugar conocido.
\end{enumerate}

\subsection*{Pistas de menor a mayor ayuda}
\begin{itemize}
\item \textbf{Si parece un callejón sin salida:} mira las paredes laterales, lee las inscripciones y revisa huecos y botones pequeños. No todos los pasos se abren con una llave.
\item \textbf{Si un mecanismo no responde:} piensa en el objeto, su forma y el lugar donde se coloca. Llevar una pieza en la mochila no equivale a usarla en el hueco correcto.
\item \textbf{Si una puerta vuelve a cerrarse:} observa si hay una placa de presión y si necesita peso. Un objeto corriente puede ser más útil en el suelo que en la mochila.
\item \textbf{Si reapareces en un sitio conocido:} puede haber un portal o un cambio de orientación. Comprueba la brújula antes de seguir contando pasos.
\item \textbf{Si encuentras objetos de piedra:} anota su forma y los lugares que tienen huecos semejantes. No los gastes ni abandones sin pensar: ayudan a conectar rutas.
\item \textbf{Si encuentras huesos o personajes amistosos:} escucha antes de atacar y conserva los restos. Algunos encuentros aportan compañeros, ayuda o pistas.
\item \textbf{Si pierdes una pelea:} vuelve con conjuros preparados, una salida libre y curación en las manos adecuadas. Evita rodearte y retírate antes de que caigan los dos de delante.
\end{itemize}

No saltes a un hueco sin preparar una posible caída. Deja reservas en un punto reconocible y anota dónde. Conserva llaves y piezas únicas; deja antes armas corrientes repetidas si falta espacio.

\clearpage
\sect{15. Dudas frecuentes y consulta rápida}
\begin{reference}{Problema}{Comprobación útil}
No puedo atacar & ¿Está el personaje consciente? ¿La mano ha dejado de estar gris? ¿Hay enemigo delante y está en primera fila para un golpe cuerpo a cuerpo?\\
El arco no dispara & Pon flechas en la otra mano y repuestos en el carcaj. Comprueba que el personaje puede usarlo.\\
No puedo ponerme un objeto & Revisa su casilla, las restricciones de clase y si llevas equipo maldito que impida cambiarlo. Una mano ocupada también limita las armas a dos manos.\\
No aparecen conjuros & Comprueba libro o símbolo, personaje, nivel de conjuro, selección en CAMP y descanso completado. Aprendido no significa preparado.\\
La magia pide mano libre & Guarda un objeto de la otra mano. Algunos conjuros crean un ataque que tendrás que usar después desde esa mano.\\
La curación no quita el veneno & Son efectos distintos: antídoto o Neutralizar veneno primero, curación de PV después.\\
No puedo descansar & Busca enemigos cercanos, resuelve el peligro y comprueba comida y estados del grupo.\\
No puedo dejar un miembro & Debe haber al menos cinco personajes y no puedes elegir al primero.\\
Un objeto mágico no hace lo esperado & Consulta su efecto y cómo usarlo en las secciones 11--13.\\
No encuentro mi partida & Comprueba el idioma elegido. Cada idioma tiene su propio guardado. Guarda la partida antes de elegir \texttt{Abandonar Juego}.\\
No se oye la aventura & La fanfarria inicial usa el sonido del MSX; los efectos de aventura y la música de la introducción requieren MoonSound. Revisa Sonidos en Preferencias.\\
\end{reference}
